%%=================================================================
%% CHAPTER 1: Secluded Majorana SIDM — Phenomenology
%% Source: github.com/0mp3s/Secluded-Majorana-SIDM (DOI: 10.5281/zenodo.19225823)
%%=================================================================

\section{Introduction: The Problem with Small Galaxies}
\label{ch1:intro}

On cosmological scales, the $\Lambda$CDM model works spectacularly well.
But zoom in to individual dwarf galaxies and something breaks: the predicted
dark matter densities in galactic cores are far higher than observed
(the ``core-cusp problem''); too many small satellite galaxies are predicted
but never seen (``missing satellites''); and the largest predicted dwarfs
are too dense to match any observed system (``too-big-to-fail'').

All three problems share a common fix: if dark matter particles scatter
off each other, they redistribute energy from the dense inner halo to
the outer region, naturally flattening the cusp into a core. This is
self-interacting dark matter (SIDM).

The question is: what mediates this self-interaction, and does it remain
consistent with everything else we know? This paper answers that question
with a single, minimal model.

\section{The Model: One Coupling, Three Parameters}
\label{ch1:model}

We consider a Majorana fermion $\chi$ (the dark matter candidate) interacting
with a real scalar mediator $\phi$ through a Yukawa interaction. The choice
of Majorana, rather than Dirac, is not arbitrary: a Majorana fermion is its
own antiparticle, which halves the relic density calculation and introduces
non-trivial quantum statistics in the self-interaction (see below).
The mediator $\phi$ has \emph{no coupling to the Standard Model} — the dark
sector is secluded. This single choice eliminates constraints from direct
detection and CMB energy injection simultaneously.

The interaction is simply:
\begin{equation}
  \Lc_{\rm int} = -\frac{y}{2}\bar{\chi}\chi\phi
\end{equation}

In the non-relativistic limit relevant to galactic dynamics, this
generates an attractive Yukawa potential between two DM particles:
\begin{equation}
  V(r) = -\frac{\alpha}{r}e^{-\mphi r}, \qquad \alpha = \frac{y^2}{4\pi}
\end{equation}
The range of the force is set by $1/m_\phi$; the strength by $\alpha$.
Three numbers — $(m_\chi,\, m_\phi,\, \alpha)$ — describe the entire model.

\subsection{Scalar Potential and Vacuum Stability}

For a stable relic $\phi$ abundance, a cubic self-coupling $\mu_3$ enables
$3\phi \leftrightarrow 2\phi$ ``cannibal'' number-changing processes that
prevent $\phi$ from overclosing the universe. This requires $\mu_3/m_\phi \gtrsim 1.7$
\cite{Farina2016}. However, a non-zero $\mu_3$ tilts the scalar potential
and can create a second vacuum deeper than $\phi = 0$, destabilizing the theory.

The full potential is:
\begin{equation}
  V(\phi) = \frac{1}{2}m_\phi^2\phi^2 + \frac{\mu_3}{3!}\phi^3 + \frac{\lambda_\phi}{4!}\phi^4
\end{equation}

Demanding that $\phi = 0$ remain the global minimum (no deeper secondary vacuum)
yields the stability condition:
\begin{equation}
  \lambda_\phi > \frac{3\mu_3^2}{8m_\phi^2}
\end{equation}
At the cannibal threshold $\mu_3/m_\phi = 1.7$, this gives $\lambda_\phi \gtrsim 1.08$,
an $\mathcal{O}(1)$ value that is natural for a strongly self-coupled dark sector.
A numerical scan confirms the tighter bound $\lambda_\phi \gtrsim 0.96$.
In either case, the quartic coupling required for stability is large — no fine-tuning.

\section{Self-Interaction Cross Section: Why a Simple ODE is Exact}
\label{ch1:vpm}

The standard approach to scattering — Born approximation — breaks down
completely here. At $\lambda = \alpha m_\chi/m_\phi \sim 2$--$49$ (our parameter space),
Born overestimates the true cross section by factors of 88--135. A full
partial-wave calculation is unavoidable.

The Variable Phase Method (VPM) solves directly for the phase shifts
$\delta_\ell$ of each partial wave via a first-order ODE:
\begin{equation}
  \frac{d\delta_\ell}{dx} = +\frac{\lambda\,e^{-x}}{\kappa\,x}
  \left[\hat{j}_\ell(\kappa x)\cos\delta_\ell - \hat{n}_\ell(\kappa x)\sin\delta_\ell\right]^2
\end{equation}
where $x = m_\phi r$ (dimensionless radius), $\kappa = \mu v_{\rm rel}/m_\phi$
(dimensionless momentum), and $\lambda = \alpha m_\chi/m_\phi$ (coupling strength).
Integrating from $x=0$ to $x \gg 1$ gives the asymptotic phase shift directly.

For identical Majorana fermions, there is a subtlety: since $\chi = \bar{\chi}$,
the two-body wavefunction must be antisymmetric under particle exchange.
This forces even-$\ell$ partial waves into the spin-singlet channel (weight 1)
and odd-$\ell$ into the spin-triplet (weight 3), giving:
\begin{equation}
  \sigmaT = \frac{2\pi}{k^2}\left[\sum_{l\,\rm even}(2l+1)\sin^2\delta_l
             + 3\sum_{l\,\rm odd}(2l+1)\sin^2\delta_l\right]
\end{equation}
The factor $2\pi$ rather than $4\pi$ accounts for identical-particle kinematics.

Why is this exact, not approximate? Chapter~2 proves that this ODE is
precisely the fluctuation determinant of the path integral over $\phi$
in the two-body background — no perturbative expansion anywhere
(see Appendix~\ref{app:vpm_proof}).

\section{Relic Density: The Same Coupling Does Both Jobs}
\label{ch1:relic}

Dark matter freezes out when the annihilation rate $\Gamma = n\langle\sigma v\rangle$
falls below the Hubble expansion rate. For Majorana $\chi\chi \to \phi\phi$
(t- and u-channel scalar exchange), the leading term is s-wave:
\begin{equation}
  \langle\sigma v\rangle_0 = \frac{\pi\alpha^2}{4m_\chi^2}
\end{equation}
This is the crucial feature of the model: the \emph{same} coupling $\alpha$
that controls the SIDM cross section at galactic velocities ($v \sim 30$--$1000$
km/s) also determines the relic abundance at freeze-out ($v \sim 0.3c$).
There is no freedom to tune one without the other.

The Sommerfeld enhancement at freeze-out is negligible: $S_0 < 1.026$
for all viable benchmarks, a correction under $2.5\%$. The secluded dark
sector means annihilation products never reach the SM plasma, so CMB
energy-injection constraints do not apply.

\section{Numerical Pipeline and Validation}
\label{ch1:pipeline}

Before presenting results, we describe the computational pipeline and its
validation, since every claim in this paper rests on the numerical cross
section solver.

\subsection{VPM Solver Implementation}
\label{ch1:solver}

The VPM ODE (Eq.~6) is integrated using a 4th-order Runge-Kutta scheme
with adaptive step size, from $x = 10^{-4}$ to $x_{\max} = 50$--$200$
(chosen to ensure $e^{-x_{\max}} < 10^{-10}$). Partial waves are summed
to $\ell_{\max}$ determined by a convergence criterion:
\begin{equation}
  \ell_{\max} = \text{smallest}\;\ell\;\text{such that}\;
  \left|\frac{\delta_\ell}{\delta_0}\right| < 10^{-6}
  \;\text{for three consecutive}\;\ell
\end{equation}
In practice, $\ell_{\max} = 3$--$8$ for BP1 ($\lambda = 1.9$) and
$\ell_{\max} = 30$--$80$ for MAP ($\lambda = 48.6$). The solver is
JIT-compiled via \texttt{numba} for performance.

\subsection{Cross-Check Suite}
\label{ch1:crosschecks}

We perform six independent validation tests on the VPM solver:

\begin{enumerate}
\item \textbf{Born limit.} At $\lambda \ll 1$, VPM must reproduce the
      analytic Born cross section. Over 5 test points with $\lambda = 0.05$--$0.75$:
      $\sigma_T^{\rm VPM}/\sigma_T^{\rm Born} = 1.000 \pm 0.002$.

\item \textbf{Unitarity.} Every phase shift must satisfy
      $|\sin^2\delta_\ell| \leq 1$. Tested on all 122 viable benchmarks
      at 9 velocities ($122 \times 9 = 1098$ evaluations): zero failures,
      worst violation $4.3\%$ below the bound.

\item \textbf{Transfer $\equiv$ elastic.}
      For identical Majorana fermions, the momentum-transfer cross section
      $\sigma_T^{\rm transfer}$ equals the elastic cross section $\sigma_T^{\rm elastic}$
      exactly, because within each spin channel the symmetrized scattering
      amplitude satisfies $|f(\theta)|^2 = |f(\pi - \theta)|^2$, making
      $\int \cos\theta\, d\sigma = 0$. We verify this identity numerically to
      $|\sigma_T^{\rm tr}/\sigma_T^{\rm el} - 1| < 10^{-12}$ at all velocities
      for both BP1 and MAP. This is a non-trivial check of the Majorana
      spin-statistics implementation.

\item \textbf{Literature comparison.}
      Validation against the Tulin, Yu \& Zurek~\cite{TYZ2013} classical-regime formula
      (their Eq.~14): agreement within $\mathcal{O}(1)$ factors in the resonant regime,
      and exact agreement in the Born regime, as expected.

\item \textbf{Blind off-grid stress test.}
      200 randomly generated parameter points never used in fitting:
      100 within the viable bounds, 50 across the full space, 50 perturbations
      of known viable points. All 200 return physical (finite, positive) cross
      sections. No solver crashes or NaN failures.

\item \textbf{Preprint consistency.} BP1 cross sections reproduced identically
      from the published code: $\sigmaT/m(30) = 0.516$,
      $\sigmaT/m(1000) = 0.072$~cm$^2$/g (0\% deviation).
\end{enumerate}

\subsection{Born Approximation Failure}
\label{ch1:born}

In our parameter space ($\lambda = 0.7$--$49$), the Born approximation
overestimates the true cross section dramatically. Table~\ref{tab:born_ratio}
shows the ratio $R_{\rm VB} = \sigma_T^{\rm VPM}/\sigma_T^{\rm Born}$ for
representative benchmarks:

\begin{table}[h!]
\centering
\caption{VPM/Born ratio $R_{\rm VB}$ at selected velocities. Born
         overestimates by factors of 2--5$\times$ at dwarf velocities,
         converging only at cluster scales ($v \gtrsim 1000$~km/s).}
\label{tab:born_ratio}
\begin{tabular}{lcccccc}
\toprule
& \multicolumn{3}{c}{BP1 ($\lambda = 1.9$)} & \multicolumn{3}{c}{MAP ($\lambda = 48.6$)} \\
\cmidrule(lr){2-4} \cmidrule(lr){5-7}
$v$ [km/s] & VPM & Born & $R_{\rm VB}$ & VPM & Born & $R_{\rm VB}$ \\
\midrule
12  & 0.430 & 1.87  & 0.23 & 1.22  & 4.84  & 0.25 \\
30  & 0.516 & 1.81  & 0.29 & 1.71  & 5.85  & 0.29 \\
200 & 0.356 & 0.742 & 0.48 & —     & —     & —    \\
1000& 0.072 & 0.085 & 0.85 & 0.203 & 0.215 & 0.95 \\
4700& 0.004 & 0.004 & 0.92 & —     & —     & —    \\
\bottomrule
\end{tabular}
\end{table}

This failure is not surprising: the Born approximation is a first-order
perturbative expansion in $\lambda$, and breaks down at $\lambda \gtrsim 1$.
What is noteworthy is the \emph{magnitude} — factors of 3--5$\times$ at the
velocities most relevant to the small-scale problems.
Any analysis of this model that relies on Born-level cross sections is unreliable.

\begin{figure}[H]
\centering
\includegraphics[width=0.75\textwidth]{fig_born_vpm}
\caption{VPM vs Born approximation and the TYZ classical-regime formula
         across the full velocity range. The Born approximation (dashed)
         overestimates by factors of 3--5$\times$ at dwarf velocities;
         VPM (solid) agrees with TYZ in the classical regime and provides
         exact results in the resonant regime.}
\label{fig:born_vpm}
\end{figure}

\subsection{Boltzmann Solver Accuracy}
\label{ch1:boltzmann}

The relic density is computed both analytically (Kolb-Turner freeze-out
approximation) and numerically (stiff ODE solver for the Boltzmann equation
$dY/dx = -\lambda_{\rm ann}[Y^2 - Y_{\rm eq}^2]$). The numerical solver
corrects the Kolb-Turner coupling by $\Delta\alpha = 16\%$ and the resulting
cross section by $\Delta(\sigmaT/m) = 19\%$ for BP1 — a correction that
exceeds many observational uncertainties and justifies the full numerical treatment.


\section{MCMC Analysis and Parameter Posteriors}
\label{ch1:mcmc}

With three free parameters and multiple observational constraints, the
critical question is whether any parameter space survives at all.

\subsection{Setup}

A Bayesian MCMC scan (\texttt{emcee}~\cite{emcee}, 32 walkers, 5000 production
steps after 500 burn-in, JIT-accelerated VPM solver) samples the posterior
over $(m_\chi, m_\phi, \alpha)$ with flat priors:
$\log_{10}(m_\chi/{\rm GeV}) \in [0.5, 2.0]$,
$\log_{10}(m_\phi/{\rm MeV}) \in [0.5, 1.5]$,
$\log_{10}\alpha \in [-4.0, -2.0]$.

The likelihood includes 13 astrophysical observables spanning four decades in velocity:
dwarf galaxy velocity dispersions ($v \sim 12$--$60$~km/s),
group-scale ellipticals ($v \sim 250$--$280$~km/s),
cluster mergers ($v \sim 1000$--$1200$~km/s),
and the Bullet Cluster ($v \sim 4700$~km/s), following
Kaplinghat, Tulin \& Yu~\cite{Kaplinghat2016}.
Upper limits (Bullet Cluster, Harvey et al.) are treated as one-sided
constraints: $\chi^2 = 0$ when $\sigmaT/m < \sigma_{\rm obs}$.

\subsection{Convergence Diagnostics}

\begin{itemize}
\item Acceptance fraction: 0.542 (within optimal range 0.2--0.7).
\item Maximum autocorrelation time: $\tau_{\max} = 75.1$ steps.
\item Effective independent samples: $N_{\rm eff} \approx 2132$.
\item Gelman-Rubin diagnostic: $\hat{R} = (1.004,\; 1.005,\; 1.005)$
      for $(m_\chi, m_\phi, \alpha)$ — well below the $\hat{R} < 1.1$ threshold.
\end{itemize}

\begin{figure}[H]
\centering
\includegraphics[width=\textwidth]{fig_chains}
\caption{MCMC chain traces for the three model parameters $(\log m_\chi,
         \log m_\phi, \log\alpha)$ across 32 walkers. Chains are well-mixed
         after the 500-step burn-in (vertical dashed line). The absence of
         long-lived modes confirms ergodicity.}
\label{fig:chains}
\end{figure}

\begin{figure}[H]
\centering
\includegraphics[width=0.65\textwidth]{fig_gelman_rubin}
\caption{Gelman-Rubin $\hat{R}$ statistic as a function of chain length.
         All three parameters converge to $\hat{R} < 1.01$ well before
         the end of the run, confirming inter-chain agreement.}
\label{fig:gelman_rubin}
\end{figure}

\begin{figure}[H]
\centering
\includegraphics[width=0.65\textwidth]{fig_autocorr}
\caption{Integrated autocorrelation time $\tau$ vs.~chain length.
         The maximum $\tau_{\max} = 75.1$ steps yields
         $N_{\rm eff} \approx 2132$ independent samples.}
\label{fig:autocorr}
\end{figure}

\subsection{Posteriors}

The parameter posteriors are unimodal with the following 68\% credible
intervals:
\begin{align}
  m_\chi &= 35.7^{+52.4}_{-25.6}\;{\rm GeV} \\
  m_\phi &= 8.4^{+4.3}_{-3.4}\;{\rm MeV} \\
  \alpha  &= 10^{-3.0 \pm 0.6}
\end{align}

A connected island of viability survives:
\begin{align}
  m_\chi &\in [10,\,100]\;{\rm GeV}, \quad
  m_\phi \in [7.6,\,14.8]\;{\rm MeV}, \quad
  \alpha \in [5\times10^{-4},\,5\times10^{-3}]
\end{align}
122 benchmark points satisfy all constraints simultaneously (80,142 raw viable,
646 representative after thinning, 394 BBN-safe). The island is narrow in
$m_\phi$ (width $\sim 7$~MeV), setting the mediator mass at a scale where
the Yukawa range $1/m_\phi \sim 25$~fm is comparable to the
de~Broglie wavelength at dwarf velocities — the resonant regime.
\begin{figure}[H]
\centering
\includegraphics[width=0.85\textwidth]{fig_corner}
\caption{Joint posterior distributions for $(\log m_\chi, \log m_\phi,
         \log\alpha)$ from the MCMC scan. Contours show 68\% and 95\%
         credible regions. The posterior is unimodal, with a strong
         $m_\chi$--$\alpha$ degeneracy driven by the relic density
         constraint.}
\label{fig:corner}
\end{figure}

\begin{figure}[H]
\centering
\includegraphics[width=0.65\textwidth]{fig_lambda_posterior}
\caption{Posterior distribution of the dimensionless coupling
         $\lambda = \alpha m_\chi/m_\phi$. The viable range
         $\lambda \in [0.7, 49]$ spans from near-Born to deeply
         non-perturbative, underscoring the necessity of full partial-wave
         calculations.}
\label{fig:lambda_posterior}
\end{figure}

\begin{figure}[H]
\centering
\includegraphics[width=0.85\textwidth]{fig_island}
\caption{Island of viability in the $(m_\chi, m_\phi)$ plane.
         Colored points are viable benchmarks; the color encodes
         $\chi^2/\nu$. The island occupies $m_\chi \in [10, 100]$~GeV,
         $m_\phi \in [7.6, 14.8]$~MeV --- a narrow band in mediator mass.}
\label{fig:island}
\end{figure}

\section{Confronting the Universe: Detailed Results}
\label{ch1:astro}

We test the model against every available astrophysical constraint,
spanning five decades in velocity from ultra-faint dwarfs ($v \sim 12$~km/s)
to the Bullet Cluster ($v \sim 4700$~km/s).

\begin{figure}[H]
\centering
\includegraphics[width=0.85\textwidth]{fig_obs_comparison}
\caption{Model predictions vs.~observational constraints across
         five decades in velocity. Points with error bars are
         astrophysical targets (Read et al., Kaplinghat et al.);
         the solid curve is the MAP best-fit. The model naturally
         transitions from strong self-interaction at dwarf scales
         to negligible interaction at cluster scales.}
\label{fig:obs_comparison}
\end{figure}

\subsection{\texorpdfstring{$\chi^2$}{Chi-squared} Goodness of Fit}

The velocity-dependent cross section is compared against 13 observational
targets from Kaplinghat et al.~\cite{Kaplinghat2016}, supplemented with
Fornax globular cluster constraints from Read et al.~\cite{Read2019}.

The unconstrained best-fit achieves:
\begin{equation}
  \chi^2_{\rm min}/\nu = 1.22/10 = 0.12
  \qquad (m_\chi = 100\,{\rm GeV},\; m_\phi = 9.2\,{\rm MeV},\;
  \alpha = 2.86\times10^{-3})
\end{equation}
Under the relic density constraint ($\Omega h^2 = 0.120 \pm 0.001$), the
best-fit is:
\begin{equation}
  \chi^2_{\rm relic}/\nu = 3.82/10 = 0.38
\end{equation}
All 17 relic-viable benchmark points achieve $\chi^2/\nu < 1$, with the
full ranking shown in Table~\ref{tab:chi2_ranking}.

\begin{figure}[H]
\centering
\includegraphics[width=0.75\textwidth]{fig_chi2_fit}
\caption{$\chi^2$ goodness-of-fit across the v34 parameter scan.
         The unconstrained minimum $\chi^2/\nu = 0.12$ drops to
         $\chi^2/\nu = 0.38$ under the relic density constraint ---
         still an excellent fit with 10 degrees of freedom.}
\label{fig:chi2_fit}
\end{figure}

\begin{table}[h!]
\centering
\caption{Best 5 relic-constrained benchmark points ranked by $\chi^2/\nu$.
         All 17 achieve $\chi^2/\nu < 1$; the full table is in
         Appendix~\ref{app:benchmarks_full}.}
\label{tab:chi2_ranking}
\begin{tabular}{cccccc}
\toprule
Rank & $m_\chi$ [GeV] & $m_\phi$ [MeV] & $\alpha$ & $\lambda$ & $\chi^2/\nu$ \\
\midrule
1 & 20.69 & 9.91  & $1.05\times10^{-3}$ & 2.1 & 0.38 \\
2 & 29.76 & 11.34 & $1.47\times10^{-3}$ & 3.9 & 0.41 \\
3 & 14.38 & 8.66  & $7.56\times10^{-4}$ & 1.3 & 0.42 \\
4 & 37.93 & 12.98 & $1.87\times10^{-3}$ & 5.5 & 0.44 \\
5 & 42.81 & 13.15 & $2.09\times10^{-3}$ & 6.8 & 0.46 \\
\bottomrule
\end{tabular}
\end{table}

\subsection{Velocity-Dependent Cross Sections}

Table~\ref{tab:sigma_vel} shows the predicted $\sigmaT/m$ at 10 characteristic
velocities alongside observational constraints. The key feature is the velocity
dependence: strong self-interaction at dwarf scales ($\sigmaT/m \sim 0.5$--$1.7$~cm$^2$/g
at $v \sim 30$~km/s) and natural suppression at cluster scales
($\sigmaT/m \sim 0.07$--$0.2$~cm$^2$/g at $v \sim 1000$~km/s).

\begin{table}[h!]
\centering
\caption{Predicted $\sigmaT/m$ [cm$^2$/g] for BP1 and MAP benchmarks
         compared with observational targets. Upper limits are one-sided.}
\label{tab:sigma_vel}
\begin{tabular}{lcccc}
\toprule
System & $v$ [km/s] & BP1 & MAP & Observed \\
\midrule
Draco           & 12   & 0.430 & 1.22  & (unconstrained) \\
TBTF dwarfs     & 30   & 0.516 & 1.71  & $1$--$3$ \\
NGC 2976        & 60   & 0.499 & —     & $\sim 0.5$ \\
NGC 720         & 250  & 0.325 & —     & $0.1$--$1$ \\
Abell 3827      & 1000 & 0.072 & 0.203 & $< 1.5$ \\
Abell 2537      & 1100 & 0.061 & —     & $< 1.0$ \\
Abell 370       & 1200 & 0.052 & —     & $< 1.0$ \\
Bullet Cluster  & 4700 & 0.004 & —     & $< 1.25$ \\
\bottomrule
\end{tabular}
\end{table}

\begin{figure}[H]
\centering
\includegraphics[width=0.85\textwidth]{fig_velocity_profile}
\caption{Self-interaction cross section $\sigmaT/m$ as a function of
         relative velocity for BP1 ($\lambda = 1.9$). The cross section
         peaks at $v \sim 30$~km/s (dwarf galaxies) and falls as $v^{-4}$
         toward cluster scales. Horizontal bands indicate observational
         targets from Kaplinghat et al.~\cite{Kaplinghat2016}.}
\label{fig:velocity_profile}
\end{figure}

\subsection{Fornax Globular Cluster Survival}
\label{ch1:fornax}

The Fornax dSph contains five globular clusters (GC1--GC5) that orbit within
$\sim 1$~kpc of the center. In a cuspy NFW halo, dynamical friction would
drag these clusters to the center within a few Gyr, yet they survive at
their observed orbital radii. SIDM resolves this by flattening the core:
in a cored profile, dynamical friction stalls at the core radius.

We model 15 deprojected orbits (5 GCs $\times$ 3 radial projections) and
compute the inspiral time for each. For MAP:
\begin{itemize}
\item 14 out of 15 deprojections stall inside the SIDM core ($r_{\rm core} = 829$~pc),
      matching the observed $r_{\rm core} = 500$--$900$~pc.
\item Jeans equation fit with Osipkov-Merritt anisotropy: $\chi^2 = 2.3$.
      For comparison, an NFW profile gives $\chi^2 = 152$.
\end{itemize}

Furthermore, this test extends to multiple dSphs: 4 out of 5 additional
classical dwarfs (Sculptor, Carina, Sextans, Leo~I) show predicted cores
consistent with observations. This is not a single object — it is a
population-level test.

\subsection{Radial Acceleration Relation}
\label{ch1:rar}

The McGaugh et al.~\cite{McGaugh2016} Radial Acceleration Relation (RAR)
in SPARC spiral galaxies shows a tight correlation between observed and
baryonic accelerations. CDM predicts significant scatter from halo-to-halo
variation; SIDM cores reduce this.

For BP1, the RAR scatter decreases from 0.177~dex (observed/CDM) to
0.122~dex, a 31\% reduction. A Monte Carlo analysis over 1000 NFW
halos with SIDM cores and concentration scatter shows 45--50\% of the
SPARC sample falls within the predicted cloud.

\subsection{Secluded Sector: Direct Detection and Cosmological Safety}
\label{ch1:secluded}

The secluded dark sector ($\phi$ has no tree-level SM coupling) has
two critical consequences:

\textbf{Direct detection.} The spin-independent cross section
$\sigma_{\rm SI} = 0$ at tree level. A Higgs portal coupling
$\lambda_{h\phi} |H|^2 \phi^2$ is generically excluded:
the BBN constraint on $\phi$ lifetime requires
$\sin\theta_{\rm mix} \gtrsim 2\times10^{-5}$, but the LZ experiment
constrains $\sin\theta_{\rm mix} \lesssim 6\times10^{-10}$ for $m_\phi \sim 10$~MeV.
The gap is a factor of $3\times10^4$ — \emph{the portal is generically closed}.

\textbf{$\Delta N_{\rm eff}$.} At BBN ($T \sim 1$~MeV),
$m_\phi/T_\phi \sim 22$, giving $\Delta N_{\rm eff} \sim e^{-22} \approx 3\times10^{-10}$
from the mediator alone.
However, the SU(2)$_d$ gauge group (Chapter~3) contributes 6 massless
bosonic degrees of freedom (3 dark gluons $\times$ 2 polarisations) that
remain deconfined at both BBN and CMB decoupling ($T_d \gg \Lambda_d$).
Entropy conservation gives
$\Delta N_{\rm eff} = \tfrac{4}{7}\!\cdot\!6\!\cdot\!\xi^4 \in [0.16, 0.26]$
for Planck-allowed decoupling temperatures,
where $\xi \equiv T_d/T_\nu$ encodes the temperature dilution from
SM entropy dumps after dark-sector decoupling.
This is a \emph{structural} prediction of the gauge group ---
not a tuneable parameter --- testable by CMB-S4
at $6$--$10\sigma$ ($\sigma_{\rm CMB\text{-}S4} = 0.027$).
Dark confinement occurs only at $z \sim 20$, well after recombination.

\textbf{Mediator relic abundance.}
$\Omega_\phi/\Omega_\chi \sim 5\times10^{-4}$ — negligible.
The cannibal mechanism ($3\phi \leftrightarrow 2\phi$) efficiently depletes
the $\phi$ population.

\subsection{Sommerfeld Enhancement at Freeze-out}

The Sommerfeld factor $S_0$ at freeze-out ($v \sim 0.3c$, $\kappa = 50$--$760$)
ranges from 1.003 to 1.025 across all 17 relic benchmarks --- a correction
of at most 2.5\%. The s-wave tree-level calculation is therefore valid.
At lower velocities ($v \sim 30$~km/s), Sommerfeld enhancement reaches
$S_0 \sim 5$--$400$, but since the dark sector is secluded, late-time
annihilation products do not inject energy into the SM plasma, and CMB
$s$-wave constraints do not apply.

\begin{figure}[H]
\centering
\includegraphics[width=0.75\textwidth]{fig_sommerfeld}
\caption{Sommerfeld enhancement factor $S_0(v)$ for representative
         benchmarks. At freeze-out ($v \sim 0.3c$), $S_0 < 1.026$.
         At galactic velocities, $S_0$ reaches $\sim 10^2$, but
         the secluded sector renders this cosmologically safe.}
\label{fig:sommerfeld}
\end{figure}

\subsection{Falsifiable Predictions}
\label{ch1:falsifiable}

The model makes several sharp, testable predictions:

\begin{enumerate}
\item \textbf{Majorana fingerprint at $v \sim 40$--$95$ km/s.}
      For MAP ($\lambda = 48.6$), the ratio $\sigma_T^{\rm Maj}/\sigma_T^{\rm Dir}$
      exceeds unity in a narrow velocity window around $v \approx 57$~km/s.
      This $R > 1$ ``bump'' is a direct consequence of odd-$\ell$ partial waves
      entering the spin-triplet channel (weight 3) and cannot occur in any
      Dirac+scalar model. Future IFU spectroscopy of intermediate dSphs
      (e.g., Leo~II, $v_{\rm max} \approx 50$~km/s) could probe this regime.

\begin{figure}[H]
\centering
\includegraphics[width=0.85\textwidth]{fig_majorana_dirac}
\caption{Majorana-to-Dirac cross-section ratio
         $R(v) = \sigma_T^{\rm Maj}/\sigma_T^{\rm Dir}$ as a function
         of relative velocity.
         \textbf{Top}: For the MAP benchmark ($\lambda = 48.6$, green),
         $R$ rises from $\sim 0.5$ at low velocity to a peak of
         $R = 1.094$ at $v \approx 53$~km/s, entering an $R > 1$
         zone (shaded) from 37 to 2000~km/s. BP1 ($\lambda = 1.9$,
         blue) remains below unity.
         \textbf{Bottom}: The logarithmic derivative $dR/d\ln v$
         shows the extremum at $v \approx 53$~km/s --- the velocity
         where intermediate dSphs should exhibit the Majorana
         fingerprint. This feature is the unique experimental
         signature of Majorana spin statistics in the self-interaction
         channel.}
\label{fig:majorana_dirac}
\end{figure}

\item \textbf{No supermassive BH seeds.}
      Gravothermal core collapse could in principle form SMBH seeds.
      For all viable benchmarks, the gravothermal timescale exceeds the
      Hubble time by a factor $> 10^4$: $t_{\rm gc}/t_{\rm univ} > 10^4$.
      This is a \emph{negative} prediction — the model does not form early SMBH seeds.

\item \textbf{Read+2019 margin.}
      The MAP point passes the Read et al.~\cite{Read2019} low-velocity constraint
      ($\sigmaT/m < 1.5$~cm$^2$/g at $v \sim 12$~km/s) with a margin of only
      $+0.05$~cm$^2$/g. A $\sim 3\%$ tighter bound would exclude MAP.

\item \textbf{Ultra-faint dwarf cores.}
      MAP predicts detectable cores in 2 out of 6 ultra-faint dwarfs at current
      sensitivity, rising to observable in future surveys (LSST, Nancy Grace Roman).
      Classical dSphs: 7 out of 8 should show cores by $t = 12$~Gyr.
\end{enumerate}
